\documentclass[10pt,a4paper]{amsart}
\usepackage[margin=19mm]{geometry}
\usepackage{amsmath,amssymb,mathtools}
\usepackage[scr=rsfs,cal=euler]{mathalfa}
\usepackage{xfrac}
\usepackage[unicode,hidelinks,bookmarks=false]{hyperref}
\newcommand{\ie}{i.e.\ }
\newcommand{\cf}{cf.\ }
\newcommand{\resp}{respectively\ }
\newcommand{\Subsection}{\subsection}
\providecommand{\nobreakdash}{\nobreak\hbox{-}}
\setlength{\emergencystretch}{2em}
\multlinegap0pt
\usepackage{amssymb}
\usepackage{mathtools}
\usepackage[shortlabels]{enumitem}
\usepackage[all]{xy}
\newtheorem{lemma}{Lemma}
\title{The Higher Connectivity at Infinity of Mapping Class Groups}
\author{Michael Mihalik}
\date{Source: arXiv:2602.12191v1 (CC BY 4.0)}
\begin{document}
\maketitle

\noindent\emph{Source and licence.} The Higher Connectivity at Infinity of Mapping Class Groups. Version arXiv:2602.12191v1 (CC BY 4.0), distributed by arXiv under the Creative Commons Attribution 4.0 International licence, http://creativecommons.org/licenses/by/4.0/, as recorded in the arXiv licence metadata for this exact version and shown on the abstract page https://arxiv.org/abs/2602.12191v1. Full licence notice: \texttt{licenses/arxiv-2602.12191-CC-BY-4.0.txt}.\par\smallskip

\noindent\textit{Editorial scope.} Counted unit: the article's lemma L8 (source0.tex lines 347--351). The article's complete original proof of it is delivered with it (source0.tex lines 352--369, one \texttt{proof} environment of 18 lines). No mathematics is written, repaired or re-derived: the statement and the proof are the article's own lines, in the article's own notation, and no retained line is substituted or re-flowed.  Retained uncounted as the source-local context the proof needs: lines 192--195; lines 199--203; lines 322--331.  One declared adaptation: exactly 1 ancillary strings inside the retained spans are omitted (image-inclusion commands and, where the source repeats a label, the repeated label definitions) and no other line differs from the source; the figure environments, with their placement, captions and labels, are kept, and the package records the omitted strings in fidelity receipts bound to the affected bodies.  No retained line contains a citation, so the wrapper carries no bibliography and no bibliography entry.  The retained lines contain the article's own diagram code, which the wrapper typesets with the standard tikz packages; no image file is delivered and the package declares no asset dependency.  Editorial formatting only, in addition: an AMS \texttt{amsart} preamble that restates the article's own declarations and macros used by the retained lines, namely the environments \texttt{lemma} on separate counters, together with the standard packages those lines need.  The retained environments print in this document as Lemma 1 in order of appearance, whereas the article prints the same items with the numbering of its own full document.  The complete native source of the article as distributed in the arXiv e-print for this exact version is bundled unchanged as \texttt{sources/194529/source0.tex}.
\begin{lemma}\label{L2}
Suppose $\{x,y\}\subset S$ such that $\{x,y\}$ generates a subgroup of $G$ isomorphic to $\mathbb Z^2$. If  $v\in \Gamma-St^{M_1(m)}(\ast)$ and 
$(v\cdot r_y)\cap St^m(\ast)\ne\emptyset$, then  $(v\cdot P(y^-,x))\cap St^m(\ast)=\emptyset$. In particular, $(vy^n\cdot l_x)\cap St^m(\ast)=\emptyset$ for $n\leq 0$ and $v\cdot r_{y^{-1}}\cap St^m(\ast)=\emptyset$. 
\end{lemma}

\begin{lemma}\label{L3}
Suppose $\{x,y,z\}\subset S^{\pm1}$ such that $\{x,y,z\}$ generates a subgroup of $G$ isomorphic to $\mathbb Z^3$. Given an integer $m>0$ there is an integer $M_2(m)>M_1(m)$ such that if  $v\in \Gamma-St^{M_2(m)}(\ast)$ and 
$(v\cdot r_x)\cap St^m(\ast)\ne\emptyset$, then  
$$(v\cdot R(x^-,y,z))\cap St^m(\ast)=\emptyset.$$ 
\end{lemma}

\begin{figure}
\vbox to 3in{\vspace {-1in} \hspace {-.1in}
\hspace{-1 in}

\vss }
\vspace{.2in}
\caption{Attaching and Mapping $Q|_{K_e}$} 
\label{Fig2}
\vspace{-.1in}
\end{figure}

\begin{lemma}\label{L8}
Suppose $x$ and $y$ are assigned respectively, to 2-cells $C_1$ and $C_2$ of $V$,  $x\ne y^{\pm1}$ and  $e=[v,w]$ is a common edge of $C_1$ and $C_2$. Suppose further that $z\in T-\{x^{\pm 1},y^{\pm 1}\}$, $v\in \partial V$ and one of the  three lines $Q(v)\cdot l_x, Q(v)\cdot l_y, Q(v)\cdot l_z$ intersects $St^{M_1(m)}(\ast)$ non-trivially then
$$Im[Q|_{K_e}([0,N]\times [0,N]\times\{0\})]\subset X-St^{M_1(m)}(\ast)(\subset X-St^m(\ast)).$$ 
(This means that the function $Q|_{K_e}$ on the front face of the cube of Figure \ref{Fig2} avoids $St^m(\ast)$. This face will be part of $S^2$, the boundary of $B^3$). 
\end{lemma}

\begin{proof} Let $u_i=v_{C_i}$ for $i\in\{1,2\}$. Suppose that $(Q(v)\cdot l_x)\cap St^{M_1(m)}(\ast)\ne\emptyset$. 
Then $(Q(u_1)\cdot l_x)\cap St^{M_1(m)+W}(\ast)\ne \emptyset$. Since $x$ was assigned to $C_1$, ($\dagger$) implies $(Q(u_1)\cdot r_{x^{-1}})\cap St^{M_1(m)+W}(\ast)\ne \emptyset$. Since $Q(v)\cdot r_{x^{-1}}$ and $Q(u_1)\cdot r_{x^{-1}}$ are parallel: 
$$(Q(v)\cdot r_{x^{-1}})\cap St^{M_1(m)+2W}(\ast)\ne \emptyset.$$ 
Since $v\in \partial V$, 
$$Q(v)\in \alpha\subset X-St^{M_2(M_1(m)+W)}(\ast)\subset X-St^{M_1(M_1(m)+2W)}(\ast).$$ 
%As $M_1(m)>m+W$, $(Q(v)\cdot r_{x^{-1}})\cap St^{M_1(m)}(\ast)\ne \emptyset$. As $M_2(m)>M_1(m)$ and $Q(v)\in X-St^{M_2(M_1(m))}(\ast)$, 
Lemma \ref{L2} (with $(M_1(m)+2W$ playing the role of $m$, $x$ playing the role of $y^{-1}$ and $y$ playing the role of $x$) implies that 
$$Q(v)\cdot P(x^+,y)\cap St^{M_1(m)+2W}(\ast)=\emptyset.$$ 
Since  $Im[Q|_{K_e}([0,N]\times [0,N]\times\{0\})]\subset Q(v)\cdot P(x^+,y)$ (see Figure \ref{Fig2}) the case when $(Q(v)\cdot l_x)\cap St^{M_1(m)}(\ast)\ne\emptyset$ is complete. 

If $(Q(v)\cdot l_y)\cap St^{M_1}(m)(\ast)\ne \emptyset$ the argument is the same (interchange  $C_1$, $u_1$, $x$ and $y$ by $C_2$, $u_2$, $y$ and $x$ respectively). 
%Next assume that $(Q(u_1)\cdot l_x)\cap St^{m+15}(\ast)= \emptyset=(Q(u_2)\cdot l_y)\cap St^{m+15}(\ast)$. 

%Next assume that $(Q(v)\cdot l_x)\cap St^{M_1(m)}(\ast)= \emptyset=(Q(v)\cdot l_y)\cap St^{M_1(m)}(\ast)$. If $(Q(v)\cdot l_z)\cap St^{M_1(m)}(\ast)\ne\emptyset$ then 
Next assume $(Q(v)\cdot r_z)\cap St^{M_1(m)}(\ast)\ne \emptyset$. Recall, $v\in X-St^{M_2(M_1(m))}(\ast)$. Lemma \ref{L3} implies that $(Q(v)\cdot R(z^-,x,y)\cap St^{M_1(m)}(\ast)=\emptyset$. Since the vertex set of $Im[Q|_{K_e}([0,N]\times [0,N]\times\{0\})]$ is a subset of $Q(v)\cdot \langle x,y\rangle\subset 
(Q(v)\cdot R(z^-,x,y)$, this case is complete. 
%Finally assume that $(Q(v)\cdot l_t)\cap St^m(\ast)=\emptyset$ for all $t\in \{x,y,z\}$. 
\end{proof}

\end{document}
